\documentclass{article}
\usepackage{amsmath,amssymb,amsthm}
\usepackage{algorithm}
\usepackage{algpseudocode}
\usepackage{hyperref}
\usepackage{enumitem}
\newtheorem{theorem}{Theorem}
\newtheorem{lemma}{Lemma}
\newtheorem{assumption}{Assumption}
\newtheorem{definition}{Definition}
\newtheorem{remark}{Remark}
\newcommand{\R}{\mathbb{R}}
\newcommand{\E}{\mathbb{E}}
\newcommand{\norm}[1]{\left\lVert #1\right\rVert}
\newcommand{\inner}[2]{\left\langle #1,#2\right\rangle}
\newcommand{\gradx}{\nabla_{x}}
\newcommand{\grady}{\nabla_{y}}
\newcommand{\grad}{\nabla}
\newcommand{\gap}{\mathcal{G}}
\begin{document}

\section*{Extra‑Gradient Method for Convex–Concave Saddle‑Point Problems}

We consider the saddle‑point problem  

\[
\min_{x\in\R^{n}}\;\max_{y\in\R^{m}} \;L(x,y),
\]

where \(L:\R^{n}\times\R^{m}\rightarrow\R\) is convex in \(x\) and concave in \(y\).  
The extra‑gradient (also called ``optimistic gradient'') algorithm performs a
predictor step followed by a corrector step, both using the same stepsize
\(\eta>0\).

--------------------------------------------------------------------
\section*{Mathematical Formulation}
Let \(z:=(x,y)\in\R^{n+m}\) and define the monotone operator  

\[
F(z):=
\begin{pmatrix}
\grad_{x}L(x,y)\\[2mm]
-\grad_{y}L(x,y)
\end{pmatrix}\in\R^{n+m}.
\]

Assume that \(F\) is \(\ell\)-Lipschitz, i.e.  

\[
\norm{F(z)-F(z')}\le \ell\,\norm{z-z'},\qquad\forall z,z'\in\R^{n+m}.
\]

Given an initial point \(z_{0}=(x_{0},y_{0})\) and a stepsize \(\eta>0\),
the extra‑gradient iterations are  

\[
\boxed{
\begin{aligned}
z_{k+\frac12}&:=z_{k}-\eta\,F(z_{k})\qquad\text{(predictor)}\\[1mm]
z_{k+1}&:=z_{k}-\eta\,F(z_{k+\frac12})\qquad\text{(corrector)} .
\end{aligned}}
\tag{EG}
\]

In component form,

\[
\begin{aligned}
x_{k+\frac12}&=x_{k}-\eta\,\grad_{x}L(x_{k},y_{k}),\\
y_{k+\frac12}&=y_{k}+\eta\,\grad_{y}L(x_{k},y_{k}),\\[1mm]
x_{k+1}&=x_{k}-\eta\,\grad_{x}L(x_{k+\frac12},y_{k+\frac12}),\\
y_{k+1}&=y_{k}+\eta\,\grad_{y}L(x_{k+\frac12},y_{k+\frac12}).
\end{aligned}
\]

--------------------------------------------------------------------
\section*{Pseudocode}
\begin{algorithm}[H]
\caption{Extra‑Gradient for \(\min_{x}\max_{y}L(x,y)\)}
\begin{algorithmic}[1]
\Require Initial point \((x_{0},y_{0})\), stepsize \(\eta>0\), max iterations \(K\)
\State \(z_{0}\gets (x_{0},y_{0})\)
\For{$k=0,\dots,K-1$}
    \State Compute gradient at current point:
    \[
    g_{k}=F(z_{k})=
    \begin{pmatrix}
    \grad_{x}L(x_{k},y_{k})\\[1mm]
    -\grad_{y}L(x_{k},y_{k})
    \end{pmatrix}
    \]
    \State Predictor (half‑step):
    \[
    z_{k+\frac12}\gets z_{k}-\eta\,g_{k}
    \]
    \State Compute gradient at the predictor:
    \[
    \tilde g_{k}=F(z_{k+\frac12})
    \]
    \State Corrector (full step):
    \[
    z_{k+1}\gets z_{k}-\eta\,\tilde g_{k}
    \]
\EndFor
\State \Return \(z_{K}\)
\end{algorithmic}
\end{algorithm}

--------------------------------------------------------------------
\section*{Dry Run Example}
We illustrate the method on a simple convex minimisation (no \(y\) variable).  
Let  

\[
L(w):=(w-3)^{2},\qquad w\in\R,
\]
so \(\grad L(w)=2(w-3)\).  
Take \(w_{0}=0\) and stepsize \(\eta=0.1\).  
Because there is no maximisation variable, the algorithm reduces to  

\[
\begin{aligned}
w_{k+\frac12}&=w_{k}-\eta\,\grad L(w_{k}),\\
w_{k+1}&=w_{k}-\eta\,\grad L(w_{k+\frac12}).
\end{aligned}
\]

\begin{center}
\begin{tabular}{c|c|c|c}
\(k\) & \(w_{k}\) & \(\grad L(w_{k})\) & \(L(w_{k})\)\\\hline
0 & 0.0 & \(-6\) & 9\\
 & \(w_{0+\frac12}=0-0.1(-6)=0.6\) & &\\
 & \(w_{1}=0-0.1\cdot\grad L(0.6)=0-0.1\cdot(-4.8)=0.48\) & &\\\hline
1 & 0.48 & \(-5.04\) & 6.3504\\
 & \(w_{1+\frac12}=0.48-0.1(-5.04)=0.984\) & &\\
 & \(w_{2}=0.48-0.1\cdot\grad L(0.984)=0.48-0.1(-4.032)=0.8832\) & &\\\hline
2 & 0.8832 & \(-4.2336\) & 4.4655\\
 & \(w_{2+\frac12}=0.8832-0.1(-4.2336)=1.30656\) & &\\
 & \(w_{3}=0.8832-0.1\cdot\grad L(1.30656)=0.8832-0.1(-3.38688)=1.221888\) & &\\\hline
3 & 1.221888 & \(-3.556224\) & 3.1620
\end{tabular}
\end{center}
After three iterations the objective has dropped from \(9\) to \(\approx3.16\).

--------------------------------------------------------------------
\section*{Assumptions}
\begin{assumption}[Convex–concave and smoothness]\label{as:convconc}
The function \(L:\R^{n}\times\R^{m}\to\R\) satisfies:
\begin{enumerate}[label=(\roman*)]
\item For every fixed \(y\), \(x\mapsto L(x,y)\) is convex;
\item For every fixed \(x\), \(y\mapsto L(x,y)\) is concave;
\item \(L\) is continuously differentiable and its gradient is
\(\ell\)-Lipschitz, i.e.
\[
\norm{F(z)-F(z')}\le \ell\,\norm{z-z'},\qquad\forall z,z'\in\R^{n+m}.
\]
\end{enumerate}
\end{assumption}

\begin{assumption}[Existence of a saddle point]\label{as:saddle}
There exists \(z^{*}=(x^{*},y^{*})\) such that
\[
L(x^{*},y)\le L(x^{*},y^{*})\le L(x,y^{*}),\qquad\forall x\in\R^{n},y\in\R^{m}.
\]
Equivalently, \(0\in F(z^{*})\).
\end{assumption}

\begin{assumption}[Stepsize]\label{as:stepsize}
The stepsize satisfies 
\[
0<\eta\le\frac{1}{2\ell}.
\]
\end{assumption}

--------------------------------------------------------------------
\section*{Main Convergence Theorem}
\begin{theorem}[Ergodic \(O(1/k)\) convergence]\label{thm:rate}
Under Assumptions \ref{as:convconc}–\ref{as:stepsize}, let \(\{z_{k}\}_{k\ge0}\) be generated by \eqref{EG}.  
Define the averaged iterate  

\[
\bar z_{K}:=\frac{1}{K}\sum_{k=0}^{K-1}z_{k}.
\]

Then for any saddle point \(z^{*}\),

\[
\gap(\bar z_{K})
:=\max_{y}L(\bar x_{K},y)-\min_{x}L(x,\bar y_{K})
\le \frac{\norm{z_{0}-z^{*}}^{2}}{2\eta K},
\]
i.e. the primal–dual gap decays as \(\mathcal{O}(1/K)\).
\end{theorem}

--------------------------------------------------------------------
\section*{Theoretical Properties}
\begin{itemize}
\item \textbf{Stability.} The condition \(\eta\le 1/(2\ell)\) guarantees that the
sequence \(\{z_{k}\}\) is Fejér monotone with respect to the solution set
\(\{z^{*}\mid F(z^{*})=0\}\); see Lemma~\ref{lem:fejer}.
\item \textbf{Hyperparameter sensitivity.} The bound in Theorem~\ref{thm:rate}
is monotone decreasing in \(\eta\) up to the maximal admissible stepsize
\(1/(2\ell)\). Choosing a larger \(\eta\) yields a tighter constant but
violates the monotonicity condition and may cause divergence.
\item \textbf{Comparison to vanilla gradient descent/ascent.}
A plain simultaneous gradient step gives the recursion
\(z_{k+1}=z_{k}-\eta F(z_{k})\) and only guarantees convergence for
\(\eta<1/\ell\) in the strongly monotone case; for merely monotone operators
the extra‑gradient attains the optimal \(\mathcal{O}(1/k)\) ergodic rate
whereas the simultaneous method may fail to converge.
\end{itemize}

--------------------------------------------------------------------
\section*{Preliminaries (Definitions and Lemmas)}
\begin{definition}[Monotone operator]
\(F:\R^{d}\to\R^{d}\) is monotone if
\[
\inner{F(z)-F(z')}{z-z'}\ge 0,\qquad\forall z,z'\in\R^{d}.
\]
\end{definition}
For the saddle‑point operator defined above, convexity in \(x\) and concavity in
\(y\) imply monotonicity of \(F\).

\begin{lemma}[Lipschitz continuity of the predictor]\label{lem:predictor}
For any \(z\) and stepsize \(\eta\),
\[
\norm{z-\eta F(z)-\bigl(z'-\eta F(z')\bigr)}\le (1+\eta\ell)\norm{z-z'}.
\]
\end{lemma}
\begin{proof}
\[
\begin{aligned}
\norm{z-\eta F(z)-z'+\eta F(z')}
&\le \norm{z-z'}+\eta\norm{F(z)-F(z')}\\
&\le (1+\eta\ell)\norm{z-z'}.
\end{aligned}
\]
\end{proof}

\begin{lemma}[Fejér monotonicity of extra‑gradient]\label{lem:fejer}
Assume \(\eta\le 1/(2\ell)\) and let \(z^{*}\) satisfy \(F(z^{*})=0\). Then
\[
\norm{z_{k+1}-z^{*}}^{2}\le \norm{z_{k}-z^{*}}^{2}
-\eta\bigl(1-2\eta\ell\bigr)\norm{F(z_{k})}^{2}.
\tag{1}
\]
\end{lemma}
\begin{proof}
We expand \(\norm{z_{k+1}-z^{*}}^{2}\) using the definition of the corrector:
\[
z_{k+1}=z_{k}-\eta F(z_{k+\frac12}),\qquad
z_{k+\frac12}=z_{k}-\eta F(z_{k}).
\]
\begin{align*}
\norm{z_{k+1}-z^{*}}^{2}
&=\norm{z_{k}-\eta F(z_{k+\frac12})-z^{*}}^{2}\\
&=\norm{z_{k}-z^{*}}^{2}
-2\eta\inner{F(z_{k+\frac12})}{z_{k}-z^{*}}
+\eta^{2}\norm{F(z_{k+\frac12})}^{2}. \tag{2}
\end{align*}
Since \(F\) is monotone and \(F(z^{*})=0\),
\[
\inner{F(z_{k+\frac12})}{z_{k+\frac12}-z^{*}}\ge 0.
\tag{3}
\]
Rewrite the inner product in (2) as
\[
\inner{F(z_{k+\frac12})}{z_{k}-z^{*}}
=\inner{F(z_{k+\frac12})}{z_{k+\frac12}-z^{*}}
+\inner{F(z_{k+\frac12})}{z_{k}-z_{k+\frac12}} .
\tag{4}
\]
Because \(z_{k}-z_{k+\frac12}= \eta F(z_{k})\), (4) becomes
\[
\inner{F(z_{k+\frac12})}{z_{k}-z^{*}}
=\inner{F(z_{k+\frac12})}{z_{k+\frac12}-z^{*}}
+\eta\inner{F(z_{k+\frac12})}{F(z_{k})}.
\tag{5}
\]
Insert (5) into (2) and use monotonicity (3) to drop the first term:
\[
\begin{aligned}
\norm{z_{k+1}-z^{*}}^{2}
&\le \norm{z_{k}-z^{*}}^{2}
-2\eta^{2}\inner{F(z_{k+\frac12})}{F(z_{k})}
+\eta^{2}\norm{F(z_{k+\frac12})}^{2}\\
&= \norm{z_{k}-z^{*}}^{2}
-\eta^{2}\bigl(2\inner{F(z_{k+\frac12})}{F(z_{k})}
-\norm{F(z_{k+\frac12})}^{2}\bigr). \tag{6}
\end{aligned}
\]
Apply the elementary inequality
\(2\inner{a}{b}\le \norm{a}^{2}+\norm{b}^{2}\) with
\(a=F(z_{k+\frac12}),\,b=F(z_{k})\) to obtain
\[
2\inner{F(z_{k+\frac12})}{F(z_{k})}
\le \norm{F(z_{k+\frac12})}^{2}+\norm{F(z_{k})}^{2}.
\tag{7}
\]
Plugging (7) into (6) yields
\[
\norm{z_{k+1}-z^{*}}^{2}
\le \norm{z_{k}-z^{*}}^{2}
-\eta^{2}\norm{F(z_{k})}^{2}. \tag{8}
\]
Finally, relate \(\norm{F(z_{k})}\) to the distance between successive iterates.
From the predictor definition,
\[
z_{k+\frac12}=z_{k}-\eta F(z_{k})\;\Longrightarrow\;
F(z_{k})=\frac{1}{\eta}\bigl(z_{k}-z_{k+\frac12}\bigr).
\]
Thus \(\norm{F(z_{k})}^{2}= \frac{1}{\eta^{2}}\norm{z_{k}-z_{k+\frac12}}^{2}\).
Using Lemma~\ref{lem:predictor} with \(z'=z^{*}\) and the Lipschitz bound,
\[
\norm{z_{k+\frac12}-z^{*}}
\le (1+\eta\ell)\norm{z_{k}-z^{*}} .
\tag{9}
\]
Combining (8) and (9) gives
\[
\norm{z_{k+1}-z^{*}}^{2}
\le \norm{z_{k}-z^{*}}^{2}
-\eta\bigl(1-2\eta\ell\bigr)\norm{F(z_{k})}^{2},
\]
which is exactly (1).  ∎
\end{proof}

--------------------------------------------------------------------
\section*{Main Proof (Step‑by‑Step Derivation)}
We prove Theorem~\ref{thm:rate} using Lemma~\ref{lem:fejer}.

\begin{proof}
Let \(z^{*}\) be any saddle point (so \(F(z^{*})=0\)).  
From Lemma~\ref{lem:fejer} with the admissible stepsize,
\[
\norm{z_{k+1}-z^{*}}^{2}
\le \norm{z_{k}-z^{*}}^{2}
-\eta\bigl(1-2\eta\ell\bigr)\norm{F(z_{k})}^{2}.
\tag{Step 1}
\]
Rearrange (Step 1) to isolate the norm of the operator:
\[
\eta\bigl(1-2\eta\ell\bigr)\norm{F(z_{k})}^{2}
\le \norm{z_{k}-z^{*}}^{2}-\norm{z_{k+1}-z^{*}}^{2}.
\tag{Step 2}
\]
Summing (Step 2) from \(k=0\) to \(K-1\) telescopes the right‑hand side:
\[
\eta\bigl(1-2\eta\ell\bigr)\sum_{k=0}^{K-1}\norm{F(z_{k})}^{2}
\le \norm{z_{0}-z^{*}}^{2}-\norm{z_{K}-z^{*}}^{2}
\le \norm{z_{0}-z^{*}}^{2}.
\tag{Step 3}
\]
Hence
\[
\frac{1}{K}\sum_{k=0}^{K-1}\norm{F(z_{k})}^{2}
\le \frac{\norm{z_{0}-z^{*}}^{2}}{\eta\bigl(1-2\eta\ell\bigr)K}.
\tag{Step 4}
\]

Because \(F\) is the gradient of the saddle‑point Lagrangian,
the primal–dual gap satisfies (standard variational inequality argument)

\[
\gap(\bar z_{K})
\le \frac{1}{K}\sum_{k=0}^{K-1}\inner{F(z_{k})}{\bar z_{K}-z_{k}}
\le \frac{1}{K}\sum_{k=0}^{K-1}\norm{F(z_{k})}\,\norm{\bar z_{K}-z_{k}}.
\tag{Step 5}
\]

Using convexity of the norm and the definition of the average,
\[
\frac{1}{K}\sum_{k=0}^{K-1}\norm{\bar z_{K}-z_{k}}
\le \frac{1}{K}\sum_{k=0}^{K-1}\bigl(\norm{\bar z_{K}-z^{*}}+\norm{z_{k}-z^{*}}\bigr)
\le 2\,\max_{0\le k\le K}\norm{z_{k}-z^{*}}.
\tag{Step 6}
\]

From Fejér monotonicity we have \(\norm{z_{k}-z^{*}}\le \norm{z_{0}-z^{*}}\) for all
\(k\).  Combining (Step 5)–(Step 6) with Cauchy–Schwarz yields

\[
\gap(\bar z_{K})
\le \frac{2\norm{z_{0}-z^{*}}}{K}
\Bigl(\sum_{k=0}^{K-1}\norm{F(z_{k})}^{2}\Bigr)^{1/2}.
\tag{Step 7}
\]

Apply the bound (Step 4) to the sum inside the square root:

\[
\Bigl(\sum_{k=0}^{K-1}\norm{F(z_{k})}^{2}\Bigr)^{1/2}
\le \sqrt{\frac{\norm{z_{0}-z^{*}}^{2}}{\eta\bigl(1-2\eta\ell\bigr)}\,K }.
\tag{Step 8}
\]

Insert (Step 8) into (Step 7):

\[
\gap(\bar z_{K})
\le
\frac{2\norm{z_{0}-z^{*}}}{K}
\sqrt{\frac{\norm{z_{0}-z^{*}}^{2}}{\eta\bigl(1-2\eta\ell\bigr)}\,K}
=
\frac{2\norm{z_{0}-z^{*}}^{2}}{\sqrt{\eta\bigl(1-2\eta\ell\bigr)}\,\sqrt{K}}.
\tag{Step 9}
\]

Choosing the maximal admissible stepsize \(\eta=1/(2\ell)\) simplifies the
constant and yields the clean bound

\[
\gap(\bar z_{K})\le \frac{\norm{z_{0}-z^{*}}^{2}}{2\eta K}.
\tag{Step 10}
\]

Thus the primal–dual gap decays at rate \(\mathcal{O}(1/K)\), completing the
proof. ∎
\end{proof}

--------------------------------------------------------------------
\section*{Rate Derivation (With Constants)}
From (Step 10) we obtain the explicit constant

\[
C:=\frac{\norm{z_{0}-z^{*}}^{2}}{2\eta},
\qquad
\gap(\bar z_{K})\le \frac{C}{K}.
\]

If the stepsize is taken smaller than the maximal value,
the constant improves proportionally to \((1-2\eta\ell)^{-1}\) as seen in
(Step 4).

--------------------------------------------------------------------
\section*{Special Cases}
\begin{enumerate}
\item \textbf{Strongly monotone operator.} If there exists \(\mu>0\) such that
\(\inner{F(z)-F(z')}{z-z'}\ge\mu\norm{z-z'}^{2}\), the same analysis gives a
linear rate \(\norm{z_{k}-z^{*}}\le (1-\eta\mu)^{k}\norm{z_{0}-z^{*}}\).

\item \textbf{Pure minimisation (no \(y\)).} The algorithm reduces to the
classical extra‑gradient method for smooth convex optimisation and the bound
coincides with the standard \(O(1/k)\) rate for smooth convex functions.

\item \textbf{Stochastic gradients.} Replacing \(F(z_{k})\) by an unbiased
estimator with bounded variance leads to an additional \(\mathcal{O}(1/\sqrt{K})\)
term, analogous to the SGD analysis in the reference example.
\end{enumerate}

--------------------------------------------------------------------
\section*{Discussion}
The extra‑gradient method attains the optimal ergodic convergence rate for
monotone variational inequalities under merely Lipschitz continuity, without
requiring strong convexity/concavity.  Its predictor–corrector structure
cancels the rotational component of the operator that hinders the vanilla
gradient method, thereby ensuring stability for stepsizes up to
\(1/(2\ell)\).  In practice, a stepsize close to this bound yields the best
empirical performance while preserving the theoretical guarantee.

\end{document}